\titleformat{\chapter}[display]{\bfseries\centering}{\huge Chapter \thechapter}{1em}{\Huge #1}
\newgeometry{top=3cm, bottom=2.2cm, left=3cm, right=2cm, headsep=1.2cm, footskip=1cm, headheight=0.3cm, textheight=24.5cm, textwidth=16cm}
\onehalfspacing
\pagestyle{fancy}
\fancyhf{}
\renewcommand{\headrulewidth}{0pt}
\cfoot{\thepage}
\chapter{Introduction}

\section{Introduction}
\noindent Social media is full of useful messages: job offers, college notices, bank alerts and chats with friends. The same channels are also used to cheat and hurt people. A fake internship may ask for a ``registration fee'' or an OTP, a link may copy a real login page, an account may be a bot, and a chat may contain threats or bullying. Most protection today comes from platform moderation and user reports, which is slow and rarely explains anything \cite{lit3}.

\noindent \textbf{Sentinel} is proposed as a response to this gap. It is a safety assistant that runs on the user's own computer. The user pastes a message, a link, a chat or an account, or switches on \emph{auto-scan} in the Chrome extension, and Sentinel tells them three things for each kind of threat: what it found, what it could not check, and what to do next. It never calls an unknown link ``safe'' or an unusual account ``fake''.

\noindent The project has two parts: a working application, and a research pipeline that trains and tests machine-learning (ML) models. A model is added to the application only if it passes a test that was fixed before training.

\section{Background Study: Terminologies and Definitions}
\begin{itemize}
\item \textbf{Phishing / Smishing:} messages that trick people into giving passwords, OTPs or money, usually through a fake website link (smishing: by SMS or chat).
\item \textbf{Social Bot / Fake Account:} an account run by a program, or one that pretends to be someone else \cite{lit3}.
\item \textbf{Harmful Language:} threats, bullying, hate and messages about self-harm. Context matters: a quoted insult is not the same as a direct one \cite{cad}.
\item \textbf{Punycode:} a way of writing non-English web addresses (\texttt{xn--}); it can make a fake site look like a real one \cite{homograph}.
\item \textbf{TF-IDF, Logistic Regression, Gradient Boosting:} standard ways of turning text into numbers and classifying it \cite{sklearn}.
\item \textbf{DistilBERT:} a small, fast language model that understands the meaning of a sentence \cite{distilbert}.
\item \textbf{Calibration:} adjusting a model's scores so that ``90\%'' really means right about 90\% of the time \cite{guo}.
\item \textbf{MinHash Fingerprint:} a short code made from a text that lets two similar messages be matched without storing the text \cite{broder}.
\item \textbf{$k$-Anonymity:} sharing something only when at least $k$ different people report it, so no single person can be identified \cite{sweeney}.
\item \textbf{Precision / Recall / FPR / AUC:} precision is how many warnings were right, recall is how many threats were caught, the false-positive rate (FPR) is how many genuine items were wrongly flagged, and AUC measures how well a model ranks bots above humans (0.5 = guessing, 1 = perfect).
\end{itemize}

\section{Fundamental Study Points}
\noindent Sentinel is built on the idea that a user should get a clear, honest explanation before acting on a suspicious message. Three fundamental ideas underpin the system:
\begin{enumerate}
\item \textbf{Separate Results, Not One Score:} a scam, a threat and a bot are different problems. Mixing them into one ``danger score'' hides the reason. Sentinel reports each one separately, with its evidence.
\item \textbf{Rules and Models Together:} rules are easy to explain but miss new wording; models catch new wording but can be wrong in surprising ways. Sentinel shows both side by side, and the model never overrides the rules.
\item \textbf{Local and Private:} all checks run on the user's own computer. The only thing ever shared is an anonymous fingerprint of a scam, and only if the user agrees.
\end{enumerate}

\section{Identification of Challenges}
\begin{itemize}
\item \textbf{Threats Come Together:} one post can contain a scam message, a fake link and a fake account at the same time.
\item \textbf{Limited Visible Data:} an ordinary user cannot see login logs or full posting histories, so many research methods \cite{lit2} cannot be used.
\item \textbf{Mixed and Disguised Language:} Indian users write in English and Hinglish, and abusers hide words (\texttt{1d10t}, \texttt{k.y.s}, \texttt{stuuupid}).
\item \textbf{Misleading High Scores:} a model can score 97\% on a biased dataset and still fail on real messages.
\item \textbf{Privacy and Platform Rules:} messages are private, and scraping profiles breaks platform terms.
\item \textbf{Limited Hardware:} the models must run quickly on a normal laptop CPU.
\end{itemize}

\section{Motivation}
\begin{itemize}
\item \textbf{Students Are Common Targets:} fake internships asking for fees, fake college or bank portals and online bullying reach students every day.
\item \textbf{Tools Without Reasons:} existing tools handle one threat at a time and give a verdict without explaining it.
\item \textbf{New Scams Spread Fast:} scammers reword the same message again and again, so each person's own filter keeps missing it.
\item \textbf{Advances in Lightweight ML:} small models such as DistilBERT can now run on a laptop, so useful checks no longer need a cloud server.
\end{itemize}

\section{Problem Statement and Proposed Solution}
\subsection{Problem Statement}
\noindent Social media users face scam messages, phishing links, harmful language and fake accounts, but current protection depends on manual moderation or separate tools that do not explain their decisions, need data the user cannot see, or send private messages to a server.

\subsection{Proposed Solution}
\noindent Sentinel is a safety assistant that runs on the user's computer, used through a four-page web app and a Chrome extension. Each request goes through separate checkers: fraud rules, harmful-language rules for English and Hinglish, a learned scam-text model, a community scam matcher, profile checks and bot detection. Every result shows the evidence found, the checks that could not be done and a next step. The extension can also scan the pages a user chooses and mark risky messages. A learned model is used only after it passes a fair test on data it has never seen.
